\section{Introduction}
\label{sec:intro}

\IEEEPARstart{P}{redicate} abstraction combined with counterexample-guided abstraction
refinement (CEGAR) has been the backbone of automated software verification for over
two decades~\cite{cegar2003,predabstraction1997,slam2002}. Tools such as
CPAChecker~\cite{cpachecker2011} implement this paradigm at industrial scale: an
abstract domain of Boolean predicates over program variables is iteratively refined
whenever model checking finds a spurious counterexample (CE). Refinement extracts
Craig interpolants from the CE's infeasibility proof~\cite{interpolation2003}, yielding
predicates that block the specific CE from recurring. This approach handles large,
realistic codebases and is a perennial top performer at SV-COMP~\cite{svcomp2022}.
Its central limitation, however, is structural: interpolation reasons about one
counterexample path in isolation. A path that unfolds a loop body $k$ times reveals
the relation between states at step $k$ and step $k+1$, but not the global invariant
that holds across all loop iterations. When the required invariant is relational---
expressing a constraint between two or more variables accumulated over the entire
loop---interpolation-based refinement produces one new predicate per unrolling and
never converges.

This limitation is best illustrated by a concrete program. Consider
\texttt{count\_up\_down-1}, a well-known benchmark from the SV-COMP Loops set. The
program runs a loop that simultaneously increments a counter \texttt{x} and decrements
a counter \texttt{y}, then asserts \texttt{y == 0} upon termination. Stock CEGAR
with interpolation finds a fresh spurious counterexample at each unrolling: first
$\{y = 0 \wedge x = n\}$, then $\{y = 1 \wedge x = n - 1\}$, then
$\{y = 2 \wedge x = n - 2\}$, and so on. Each predicate is locally valid for
one loop iteration but insufficient to prove the property globally.
After \CUDStockRefs{} refinements the tool exhausts its \CUDStockWall{}\,s budget
and returns \textsc{Unknown}. The invariant that would close the proof---$x + y = n$,
which captures the joint evolution of both counters---is never discovered because no
single CE trace makes it visible.

Large language models pre-trained on code corpora have seen thousands of programs with
analogous loop structures. Given the source code of \texttt{count\_up\_down-1}, an LLM
trained on competitive programming and systems software can plausibly identify $x + y = n$
as the relevant invariant from structural pattern matching alone. The critical challenge
is that LLM output is untrusted: an incorrect predicate injected into the abstract
domain could cause CEGAR to terminate with a wrong safety verdict, violating the
soundness guarantee that makes formal verification meaningful. Simply asking the model
for predicates and using them directly is not an option.

VGuide resolves this tension with a strict oracle principle: \emph{the LLM proposes,
SMT certifies, CEGAR decides}. At no point does an unvalidated LLM output influence
the abstract state or the verification verdict. VGuide fires exactly once per analysis,
on the first spurious CE, and channels LLM suggestions through a three-tier validation
cascade before any predicate reaches the CEGAR refinement machinery. Predicates that
fail validation are silently discarded. CEGAR retains full control of the refinement
loop; VGuide is a hint supplier, not a decision maker.

This paper makes the following contributions:
\begin{enumerate}
  \item \textbf{VGuide system.} A CE-guided LLM predicate oracle integrated into
        CPAChecker's PredicateCPA with formal soundness guarantees. VGuide intercepts
        the CEGAR loop at the first spurious CE, assembles a structured context pack,
        and injects only SMT-validated predicates at CFA loop heads.
  \item \textbf{Three-tier validation cascade.} A layered filtering architecture
        comprising contract checking (L1), scope-vocabulary filtering (L2), and Z3
        SMT entailment (L3) that maintains the standard CEGAR soundness invariant
        under arbitrary LLM outputs.
  \item \textbf{Empirical evaluation.} On 764 tasks from the SV-COMP Loops
        ReachSafety benchmark, VGuide gains \LoopsDelta{} tasks over the base
        predicate analysis (\LoopsStock{}$\to$\LoopsVGuide{}) and \PortfolioLoopsDelta{}
        tasks over the full svcomp26 competition portfolio
        (\PortfolioLoopsStock{}$\to$\PortfolioLoopsVGuide{}). In a
        competition-grade combined evaluation, the net gain is \CompTotalDelta{} tasks
        (\CompReachDelta{} ReachSafety, \CompOvfDelta{} NoOverflow), with zero wrong
        verdicts across all configurations.
  \item \textbf{Ablation.} A source-prior variant that fires the LLM at analysis
        start with source code only---no CE trace---yields zero improvement
        (\AblationSourcePrior{} tasks, identical to stock). This establishes that the
        value of VGuide comes entirely from reasoning about the CE trace, not from
        static program understanding, and confirms that CE-guided refinement is the
        correct integration point.
\end{enumerate}

The remainder of this paper is organized as follows. \Cref{sec:background} reviews
predicate abstraction and CEGAR. \Cref{sec:design} presents the VGuide architecture.
\Cref{sec:evaluation} reports evaluation results. \Cref{sec:discussion} discusses
case studies, failure modes, and limitations. \Cref{sec:related} surveys related
work. \Cref{sec:conclusion} concludes.
